\documentclass[12pt,a4paper]{report}

% ── Encodage & langue ────────────────────────────────────────────────
\usepackage[utf8]{inputenc}
\usepackage[T1]{fontenc}
\usepackage[french]{babel}

% ── Mise en page ─────────────────────────────────────────────────────
\usepackage[top=2.5cm, bottom=2.5cm, left=2.5cm, right=2.5cm]{geometry}
\usepackage{setspace}
\onehalfspacing
\setlength{\parindent}{1em}
\setlength{\parskip}{0.4em}

% ── Mathématiques ────────────────────────────────────────────────────
\usepackage{amsmath, amssymb}

% ── Tableaux ─────────────────────────────────────────────────────────
\usepackage{booktabs}
\usepackage{array}
\usepackage{multirow}
\usepackage{longtable}
\usepackage{tabularx}
\usepackage{colortbl}

% ── Graphiques & couleurs ────────────────────────────────────────────
\usepackage{graphicx}
\usepackage[dvipsnames,table]{xcolor}
\usepackage{tikz}
\usetikzlibrary{arrows.meta, positioning, shapes.geometric, fit, calc, backgrounds, decorations.pathreplacing, patterns, shadows}
\usepackage{float}

% ── Listings / pseudocode ────────────────────────────────────────────
\usepackage{listings}
\lstset{
  basicstyle=\ttfamily\footnotesize,
  keywordstyle=\bfseries\color{MidnightBlue},
  commentstyle=\itshape\color{Gray},
  stringstyle=\color{OliveGreen},
  frame=single,
  framerule=0.4pt,
  breaklines=true,
  tabsize=2,
  xleftmargin=1em,
  numbers=left,
  numberstyle=\tiny\color{Gray},
  numbersep=8pt,
  captionpos=b,
  backgroundcolor=\color{gray!5},
  rulecolor=\color{gray!40},
}

% ── Hyperliens ───────────────────────────────────────────────────────
\usepackage[hidelinks]{hyperref}

% ── En-têtes / pieds de page ────────────────────────────────────────
\usepackage{fancyhdr}
\pagestyle{fancy}
\fancyhf{}
\fancyhead[L]{\small\itshape Chapitre V --- Développement Logiciel}
\fancyhead[R]{\small\itshape PFE : CNC 5 axes}
\fancyfoot[C]{\thepage}
\renewcommand{\headrulewidth}{0.4pt}

% ── Captions ─────────────────────────────────────────────────────────
\usepackage{caption}
\usepackage{subcaption}
\captionsetup{labelfont=bf, font=small, skip=6pt}

% ── Encadrés ─────────────────────────────────────────────────────────
\usepackage{tcolorbox}
\tcbuselibrary{skins, breakable}

% ── Commandes personnalisées ─────────────────────────────────────────
\newcommand{\code}[1]{\texttt{\small #1}}
\newcommand{\screenshot}[2]{%
  \begin{figure}[H]
  \centering
  \fbox{\begin{minipage}{0.85\textwidth}
  \centering
  \vspace{3cm}
  {\Large\color{gray!60} \textbf{CAPTURE D'ÉCRAN}}\\[6pt]
  {\large\color{gray!80} #1}\\[4pt]
  {\small\color{gray!50}\itshape Remplacer par : \texttt{figures/#2}}
  \vspace{3cm}
  \end{minipage}}
  \caption{#1}
  \label{fig:#2}
  \end{figure}
}

% ══════════════════════════════════════════════════════════════════════
\begin{document}

\setcounter{chapter}{4}

\chapter{Développement Logiciel et Ingénierie de Commande}

% ══════════════════════════════════════════════════════════════════════
\section{Introduction}
% ══════════════════════════════════════════════════════════════════════

Le pilotage d'une machine CNC 5~axes ne se résume pas à la génération d'impulsions Step/Dir par un firmware embarqué. Si le contrôleur FluidNC, exécuté sur l'ESP32 (cf.~Chapitre~IV), assure la conversion temps réel du G-code en signaux moteur, il ne constitue qu'un maillon d'une chaîne logicielle plus vaste. La supervision de l'état machine en temps réel, la résolution de la cinématique inverse propre à l'architecture Trunnion, la planification anticipative des trajectoires, la protection active des efforts de coupe et l'interface de commande opérateur nécessitent un logiciel de bord dédié --- l'application \textbf{Forgeron}.

Ce chapitre présente le développement logiciel selon cinq piliers~:
\begin{itemize}
  \item L'\textbf{architecture logicielle} (section~V.2) --- structuration en Clean Architecture, modélisation UML.
  \item Les \textbf{algorithmes de cinématique inverse} (section~V.3) --- modèle géométrique Trunnion, RTCP, singularités.
  \item La \textbf{planification de trajectoire} et le streaming industriel (section~V.4).
  \item Le développement de l'\textbf{IHM de supervision} multi-plateforme (section~V.5).
  \item La \textbf{sûreté de fonctionnement} du code embarqué et applicatif (section~V.6).
\end{itemize}

% ══════════════════════════════════════════════════════════════════════
\section{Architecture Logicielle et Modélisation UML}
% ══════════════════════════════════════════════════════════════════════

\subsection{Philosophie « Clean Architecture »}

L'application Forgeron est structurée selon le paradigme \textit{Clean Architecture} de Robert C.~Martin~\cite{martin2017clean}. Ce modèle impose une séparation stricte en couches concentriques, chacune ayant un niveau d'abstraction croissant vers le centre. La règle fondamentale --- la \textbf{Dependency Rule} --- stipule que les dépendances ne pointent que vers l'intérieur~: les couches internes ne connaissent jamais les couches externes.

La figure~\ref{fig:clean_arch} illustre l'organisation concentrique des quatre couches du projet.

\begin{figure}[H]
\centering
\begin{tikzpicture}[
  layer/.style={draw, rounded corners=8pt, thick, align=center, font=\small},
]
  % Couche externe → interne
  \node[layer, fill=blue!8, minimum width=14cm, minimum height=9cm] (pres) {};
  \node[layer, fill=green!8, minimum width=10.5cm, minimum height=6.5cm] (data) {};
  \node[layer, fill=orange!10, minimum width=7cm, minimum height=4cm] (app) {};
  \node[layer, fill=red!8, minimum width=3.8cm, minimum height=2cm] (dom) {};

  % Labels
  \node[font=\bfseries\small, anchor=north west] at ([shift={(0.3,-0.3)}]pres.north west) {\textcolor{blue!70}{PRESENTATION}};
  \node[font=\bfseries\small, anchor=north west] at ([shift={(0.3,-0.3)}]data.north west) {\textcolor{green!50!black}{DATA}};
  \node[font=\bfseries\small, anchor=north west] at ([shift={(0.3,-0.3)}]app.north west) {\textcolor{orange!70!black}{APPLICATION}};
  \node[font=\bfseries\small] at (dom.center) {\textcolor{red!60!black}{DOMAIN}};

  % Contenu Presentation
  \node[font=\scriptsize\ttfamily, anchor=south west] at ([shift={(0.4,0.4)}]pres.south west) {MainScaffold, DashboardScreen, MobileProbingScreen};
  \node[font=\scriptsize\ttfamily, anchor=south east] at ([shift={(-0.4,0.4)}]pres.south east) {widgets/, tutorial/, AppColors};

  % Contenu Data
  \node[font=\scriptsize\ttfamily, anchor=south west] at ([shift={(0.3,0.3)}]data.south west) {FluidNCMachineRepository, GrblParser, FluidNCConnection};
  \node[font=\scriptsize\ttfamily, anchor=south east] at ([shift={(-0.3,0.3)}]data.south east) {MockMachineRepository, FluidNCHttpClient};

  % Contenu Application
  \node[font=\scriptsize\ttfamily, anchor=south] at ([shift={(0,-0.1)}]app.south) {ForceGuardService, GCodeStreamingService, machineRepositoryProvider};

  % Contenu Domain
  \node[font=\scriptsize\ttfamily, anchor=south] at ([shift={(0,0.15)}]dom.south) {MachineState, MachineRepository};
  \node[font=\scriptsize\ttfamily, anchor=north] at ([shift={(0,-0.15)}]dom.north) {TrunnionConfig, MachiningMode};

  % Flèches Dependency Rule
  \draw[-{Stealth[length=3mm]}, very thick, dashed, gray!60] ([shift={(0.5,0)}]pres.east) -- ++(1.2,0) node[right, font=\scriptsize\itshape, text=gray]{Dependency Rule} -- ++(0,-3.5) -- ++(-1.2,0);

\end{tikzpicture}
\caption{Architecture concentrique Clean Architecture du projet Forgeron.}
\label{fig:clean_arch}
\end{figure}

Le tableau~\ref{tab:clean_arch} détaille le contenu de chaque couche.

\begin{table}[H]
\centering
\caption{Mapping Clean Architecture du projet Forgeron.}
\label{tab:clean_arch}
\begin{tabularx}{\textwidth}{>{\bfseries}l p{3cm} X}
\toprule
Couche & Rôle & Fichiers et classes clés \\
\midrule
\rowcolor{red!5}
Domain & Entités immuables + contrats (interfaces abstraites) & \code{MachineState} (Freezed, 9~états FSM), \code{TrunnionConfig} (paramètres D-H), \code{MachiningMode} (enum 3AX/5AX), \code{Tool}, \code{Alarm}, \code{WCS}, \code{Macro} --- \code{MachineRepository} (abstract, 14~méthodes), \code{ConfigRepository}, \code{FileRepository} \\
\rowcolor{orange!5}
Application & Orchestration, injection de dépendances, services métier & \code{machineRepositoryProvider} (DI Riverpod), \code{machineStateProvider} (StreamProvider), \code{ForceGuardService} (bridage efforts), \code{GCodeStreamingService} (character-counting), \code{AutoDiscoveryService} \\
\rowcolor{green!5}
Data & Implémentations concrètes des contrats & \code{FluidNCMachineRepository} (WebSocket), \code{GrblParser} (289~lignes), \code{FluidNCConnection} (backoff + heartbeat), \code{MockMachineRepository} (simulation) \\
\rowcolor{blue!5}
Presentation & Interface utilisateur Flutter & \code{MainScaffold} (responsive), 6~écrans desktop + 4~écrans mobile, widgets industriels (\code{GlassPanel}, DRO, jauges) \\
\bottomrule
\end{tabularx}
\end{table}

Cette architecture garantit un avantage déterminant~: la substitution de l'implémentation réelle par un simulateur se fait par \textbf{une seule ligne} dans le provider d'injection de dépendances~:

\begin{lstlisting}[language=Java, caption={Provider DI --- bascule Mock/FluidNC (inversion de dépendance).}, label={lst:di}]
final machineRepositoryProvider = Provider<MachineRepository>((ref) {
  final isSim = ref.watch(isSimulationModeProvider);
  if (isSim) return MockMachineRepository();    // <-- Simulateur
  final conn = FluidNCConnection(wsUrl);         // <-- ESP32 reel
  return FluidNCMachineRepository(conn);
});
\end{lstlisting}

\subsection{Modélisation structurelle et comportementale}

\subsubsection{Machine à états finis (FSM)}

La classe \code{MachineState}, générée par Freezed (immuable, sérialisable), modélise l'état complet de la machine selon neuf états conformes au protocole GRBL~v1.1. La figure~\ref{fig:fsm} présente le diagramme de transitions.

\begin{figure}[H]
\centering
\begin{tikzpicture}[
  state/.style={rectangle, draw, rounded corners=4pt, minimum width=2cm, minimum height=0.8cm, font=\small\ttfamily\bfseries, fill=#1!15, draw=#1!60},
  trans/.style={-{Stealth[length=2.5mm]}, thick, color=gray!70},
  label/.style={font=\tiny\itshape, text=gray!80, fill=white, inner sep=1pt},
  node distance=1.8cm and 2.2cm,
]
  % États
  \node[state=gray]   (off)   {OFFLINE};
  \node[state=green, right=of off]  (idle)  {IDLE};
  \node[state=blue, right=of idle]  (run)   {RUN};
  \node[state=orange, below=of run] (hold)  {HOLD};
  \node[state=red, below=of idle]   (alarm) {ALARM};
  \node[state=cyan, right=of run]   (home)  {HOME};
  \node[state=purple, below=of home](check) {CHECK};
  \node[state=brown, right=of home] (door)  {DOOR};
  \node[state=gray!60, below=of door](sleep){SLEEP};

  % Transitions principales
  \draw[trans] (off)   -- node[label, above]{connect}   (idle);
  \draw[trans] (idle)  -- node[label, above]{cycle start} (run);
  \draw[trans] (run)   -- node[label, right]{feed hold}  (hold);
  \draw[trans] (hold)  -- node[label, below]{resume \textasciitilde} (run);
  \draw[trans] (hold.west)  -- node[label, below]{reset}  (idle.south);
  \draw[trans] (run.south west)   -- node[label, left, pos=0.3]{ALARM:N}   (alarm);
  \draw[trans] (alarm) -- node[label, left]{\$X unlock} (idle);
  \draw[trans] (idle)  -- node[label, above]{\$H}       (home);
  \draw[trans] (home)  -- node[label, right]{done}      (idle.north east);
  \draw[trans] (idle.east) to[bend left=15] node[label, above right]{\$C} (check.west);
  \draw[trans] (check) -- node[label, right]{\$C off}  (idle.south east);
  \draw[trans] (run)   -- node[label, above]{door open} (door);
  \draw[trans] (door)  -- node[label, right]{door close}(run.east);
  \draw[trans] (idle.south east) to[bend right=40] node[label, right]{\$SLP} (sleep);
  \draw[trans] (sleep.north west) to[bend right=40] node[label, left]{wake} (idle.south east);

  % Perte connexion
  \draw[trans, dashed, red!50] (run.north) to[bend left=60] node[label, above]{timeout 10s} (off.north);
\end{tikzpicture}
\caption{Diagramme de machine à états finis --- FSM de \code{MachineState} (9~états, protocole GRBL v1.1).}
\label{fig:fsm}
\end{figure}

Les transitions sont déclenchées par les \textit{status reports} émis par FluidNC au format \code{<State|MPos:...|FS:...>}. Le parser central \code{GrblParser.parse()} dispatche chaque message entrant vers l'une des cinq méthodes spécialisées (tableau~\ref{tab:parser_methods}).

\begin{table}[H]
\centering
\caption{Méthodes de dispatch du parser \code{GrblParser}.}
\label{tab:parser_methods}
\begin{tabularx}{\textwidth}{l l X}
\toprule
Préfixe du message & Méthode appelée & Données extraites \\
\midrule
\code{<State|...>} & \code{parseStatusReport()} & Positions MPos/WPos (5~axes), feedrate, broche, overrides, fins de course, buffers \\
\code{ALARM:N} & \code{parseAlarm()} & Code numérique $\rightarrow$ transition FSM vers \code{alarm} \\
\code{[GC:...]} & \code{parseModalState()} & WCS actif (G54--G59.3), outil (T\textit{n}), avance (F), broche (S) \\
\code{[PRB:...]} & \code{parseProbeReport()} & Coordonnées de contact palpeur [X,Y,Z], succès/échec \\
\code{[MSG:...]} & \code{parseMessage()} & Messages informatifs FluidNC \\
\bottomrule
\end{tabularx}
\end{table}

\subsubsection{Diagramme de classes}

La figure~\ref{fig:class_diagram} illustre les relations structurelles entre les composants principaux de la couche Data et leur lien avec les contrats du Domain.

\begin{figure}[H]
\centering
\begin{tikzpicture}[
  class/.style={rectangle, draw, rounded corners=2pt, minimum width=4.2cm, minimum height=1.4cm, fill=blue!5, font=\small\ttfamily, align=center, drop shadow={shadow xshift=1pt, shadow yshift=-1pt, opacity=0.15}},
  abstract/.style={class, fill=orange!10, font=\small\ttfamily},
  interface/.style={class, fill=yellow!10, font=\small\itshape\ttfamily},
  arrow/.style={-{Stealth[length=3mm]}, thick},
  dasharrow/.style={-{Stealth[length=3mm]}, thick, dashed},
  note/.style={font=\tiny\itshape, text=gray},
]
  % Domain layer
  \begin{scope}[on background layer]
    \node[draw=red!30, fill=red!3, rounded corners=6pt, minimum width=10cm, minimum height=2.2cm, label={[font=\tiny\bfseries\color{red!50}]above:DOMAIN}] at (2.5,5.5) {};
  \end{scope}
  \node[interface] (repo) at (0,5.5) {$\ll$abstract$\gg$\\MachineRepository};
  \node[class, fill=red!8] (state) at (5.5,5.5) {MachineState\\{\tiny\itshape (Freezed, immuable)}};

  % Data layer
  \begin{scope}[on background layer]
    \node[draw=green!30, fill=green!3, rounded corners=6pt, minimum width=14cm, minimum height=5.5cm, label={[font=\tiny\bfseries\color{green!50!black}]above:DATA}] at (2.5,2) {};
  \end{scope}
  \node[class, fill=green!10] (fluidnc) at (-2,3) {FluidNCMachine\\Repository};
  \node[class, fill=green!10] (mock) at (3,3) {MockMachine\\Repository};
  \node[class] (conn) at (-2,0.5) {FluidNCConnection\\{\tiny WebSocket + Backoff}};
  \node[class] (parser) at (3,0.5) {GrblParser\\{\tiny 289 lignes, 5 formats}};
  \node[class] (ws) at (7,0.5) {WebSocketChannel\\{\tiny dart:io}};

  % Relations
  \draw[dasharrow] (fluidnc) -- node[left, note]{implements} (repo);
  \draw[dasharrow] (mock) -- node[right, note]{implements} (repo);
  \draw[arrow] (fluidnc) -- node[left, note]{utilise} (conn);
  \draw[arrow] (fluidnc) -- node[above, note]{parse via} (parser);
  \draw[arrow] (conn) -- node[above, note]{wraps} (ws);
  \draw[arrow, dashed, gray] (parser) -- node[right, note]{produit} (state);

\end{tikzpicture}
\caption{Diagramme de classes simplifié --- couches Domain et Data.}
\label{fig:class_diagram}
\end{figure}

\subsection{Gestion du multitâche sous FreeRTOS}

L'ESP32 dispose de deux c\oe urs matériels exploités par FreeRTOS pour garantir le déterminisme du pilotage multi-axes. La figure~\ref{fig:dual_core} illustre la répartition des tâches.

\begin{figure}[H]
\centering
\begin{tikzpicture}[
  core/.style={draw, thick, rounded corners=6pt, minimum width=6cm, minimum height=4cm, fill=#1!8},
  task/.style={draw, rounded corners=3pt, minimum width=5cm, minimum height=0.7cm, fill=#1!15, font=\small},
]
  % Core 0
  \node[core=blue, label={[font=\bfseries]above:C\oe ur 0 --- Réseau \& Parsing}] (c0) at (0,0) {};
  \node[task=blue] at (0,1.2) {Pile Wi-Fi (TCP/IP, mDNS)};
  \node[task=blue] at (0,0.3) {Serveur WebSocket (port 80)};
  \node[task=blue] at (0,-0.6) {Serveur HTTP (upload, config YAML)};
  \node[task=blue] at (0,-1.5) {Parsing GRBL + JSON};

  % Core 1
  \node[core=red, label={[font=\bfseries]above:C\oe ur 1 --- Temps Réel (ISR)}] (c1) at (8,0) {};
  \node[task=red] at (8,1.2) {Timer HW 32~kHz ($T = 31{,}25\;\mu$s)};
  \node[task=red] at (8,0.3) {ISR Step/Dir (5 axes)};
  \node[task=red] at (8,-0.6) {Planner Look-Ahead (15 blocs)};
  \node[task=red] at (8,-1.5) {S-Curve Acceleration Profile};

  % Séparation
  \draw[very thick, dashed, gray!50] (4,-2.5) -- (4,2.8);
  \node[font=\scriptsize\itshape, text=gray, rotate=90] at (4,0) {Isolation matérielle};

\end{tikzpicture}
\caption{Répartition des tâches FreeRTOS sur les deux c\oe urs de l'ESP32.}
\label{fig:dual_core}
\end{figure}

La charge CPU du c\oe ur critique (c\oe ur~1) se calcule comme suit~:

\begin{equation}
\text{Charge}_{\text{CPU}_1} = \frac{N_{axes} \times t_{ISR}}{T_{timer}} = \frac{5 \times 2\;\mu\text{s}}{31{,}25\;\mu\text{s}} = 32\%
\label{eq:cpu_load}
\end{equation}

\noindent avec $N_{axes} = 5$, $t_{ISR} \approx 2\;\mu$s par axe et $T_{timer} = 1/32\,000 = 31{,}25\;\mu$s. La marge résiduelle de \textbf{68\%} garantit l'absence de \textit{jitter} sur les signaux Step, même en cas de pic de charge réseau sur le c\oe ur~0.

% ══════════════════════════════════════════════════════════════════════
\section{Cinématique Inverse de l'Architecture Trunnion}
% ══════════════════════════════════════════════════════════════════════

\subsection{Définition du modèle géométrique}

Le Trunnion (Table-Table) comporte deux axes rotatifs superposés~: l'axe~A (berceau, rotation autour de~X, plage $\pm 90°$) et l'axe~C (plateau, rotation autour de~Z, plage $360°$). La figure~\ref{fig:trunnion_kin} illustre le modèle géométrique.

\begin{figure}[H]
\centering
\begin{tikzpicture}[scale=0.9,
  axis/.style={-{Stealth[length=3mm]}, very thick},
  dim/.style={|{Stealth[length=2mm]}-{Stealth[length=2mm]}|, thin, gray},
]
  % Base / table
  \fill[gray!10] (-3,-0.3) rectangle (3,0);
  \draw[thick] (-3,0) -- (3,0);
  \node[font=\footnotesize, below] at (0,-0.3) {Base machine};

  % Berceau (axe A)
  \draw[thick, blue!60, fill=blue!5] (-2.5,0) -- (-2.5,3) arc(180:0:2.5) -- (2.5,0);
  \node[font=\footnotesize\bfseries, blue!70] at (3.2,1.5) {Berceau};

  % Pivot A
  \fill[red!70] (0,3) circle(4pt);
  \node[font=\scriptsize\bfseries, red!70, above right] at (0.15,3.1) {Pivot A};
  \draw[axis, red!70] (0,3) -- (2,3) node[right, font=\scriptsize]{Axe A (rot. autour de X)};

  % Plateau (axe C)
  \draw[thick, green!60!black, fill=green!8] (-1.5,2.2) rectangle (1.5,2.7);
  \node[font=\footnotesize\bfseries, green!60!black] at (0,2.45) {Plateau C};
  \draw[axis, green!60!black] (0,2.45) ++(0,0) -- ++(0,1.2) node[above, font=\scriptsize]{Axe C (rot. autour de Z)};

  % Pièce
  \fill[orange!30] (-0.8,2.7) rectangle (0.8,3.3);
  \node[font=\scriptsize\bfseries] at (0,3.0) {Pièce};

  % Outil
  \draw[very thick, brown!70] (0,5.5) -- (0,3.5);
  \fill[brown!50] (-0.15,5.5) rectangle (0.15,3.5);
  \node[font=\footnotesize\bfseries, brown!70, right] at (0.2,4.5) {Outil};
  \fill[brown!70] (0,3.5) circle(2pt);
  \node[font=\scriptsize, brown!70, left] at (-0.2,3.5) {TCP};

  % Distance d
  \draw[dim] (1.8,3) -- node[right, font=\scriptsize]{$d = 45$ mm} (1.8,2.45);

  % Axes cartésiens
  \draw[axis, gray] (-3.5,0) -- (-3.5,2) node[above, font=\scriptsize]{Z};
  \draw[axis, gray] (-3.5,0) -- (-1.5,0) node[right, font=\scriptsize]{X};
  \fill[gray] (-3.5,0) circle(2pt);
  \node[font=\scriptsize, gray, below left] at (-3.5,0) {Y (vers l'obs.)};

\end{tikzpicture}
\caption{Modèle géométrique du Trunnion Table-Table --- axes A, C, distance pivot-table $d$.}
\label{fig:trunnion_kin}
\end{figure}

Les paramètres géométriques, codés dans la classe \code{TrunnionConfig}, sont résumés dans le tableau~\ref{tab:dh_params}.

\begin{table}[H]
\centering
\caption{Paramètres géométriques du Trunnion (source~: \code{TrunnionConfig}).}
\label{tab:dh_params}
\begin{tabular}{l r l}
\toprule
Paramètre & Valeur & Attribut code \\
\midrule
Plage axe A (berceau) & $\pm 90°$ & \code{aAxisMaxAngle = 90.0} \\
Plage axe C (plateau) & $360°$ continu & \code{cAxisMaxAngle = 360.0} \\
Réduction GT2 & 6:1 (120/20) & \code{drivenPulleyTeeth / drivingPulleyTeeth} \\
Offset pivot--table & 45~mm & \code{pivotToTableOffset = 45.0} \\
Micro-stepping & 1/16 & \code{microstepping = 16} \\
Pas moteur & 200~pas/tr & \code{motorStepsPerRev = 200} \\
\bottomrule
\end{tabular}
\end{table}

La résolution angulaire effective est déduite du nombre de pas par degré~:

\begin{equation}
\text{stepsPerDegree} = \frac{\text{motorStepsPerRev} \times \text{microstepping} \times \text{reductionRatio}}{360} = \frac{200 \times 16 \times 6}{360} = 53{,}33 \;\text{pas/°}
\label{eq:steps_per_deg}
\end{equation}

\begin{equation}
\boxed{\Delta\theta = \frac{1}{53{,}33} \approx 0{,}019°}
\label{eq:angular_res}
\end{equation}

Cette résolution est largement suffisante pour l'usinage d'aluminium AW-2017A ($\pm 0{,}1°$ typique).

\subsection{Algorithme de résolution analytique}

La cinématique inverse transforme les coordonnées pièce $(X_w, Y_w, Z_w)$ en coordonnées machine $(X_m, Y_m, Z_m)$ pour des angles $(A, C)$ donnés. Soit $d$ la distance pivot-table~:

\begin{align}
X_m &= X_w + d \cdot \sin(A) \cdot \sin(C) \label{eq:kin_x} \\
Y_m &= Y_w - d \cdot \sin(A) \cdot \cos(C) \label{eq:kin_y} \\
Z_m &= Z_w + d \cdot (1 - \cos(A)) \label{eq:kin_z}
\end{align}

\begin{lstlisting}[caption={Pseudocode --- cinématique inverse Trunnion.}, label={alg:kin_inv}]
function inverseKinematics(Xw, Yw, Zw, A_deg, C_deg, d=45):
    A = A_deg * PI / 180
    C = C_deg * PI / 180
    Xm = Xw + d * sin(A) * sin(C)
    Ym = Yw - d * sin(A) * cos(C)
    Zm = Zw + d * (1 - cos(A))
    return (Xm, Ym, Zm, A_deg, C_deg)
\end{lstlisting}

\paragraph{Exemple numérique.} Pour $X_w = 50$~mm, $Y_w = 30$~mm, $Z_w = 0$, $A = 30°$, $C = 45°$~:
\begin{align*}
X_m &= 50 + 45 \times \sin(30°) \times \sin(45°) = 50 + 45 \times 0{,}5 \times 0{,}707 = \mathbf{65{,}9\;\text{mm}} \\
Y_m &= 30 - 45 \times \sin(30°) \times \cos(45°) = 30 - 15{,}9 = \mathbf{14{,}1\;\text{mm}} \\
Z_m &= 0 + 45 \times (1 - \cos(30°)) = 45 \times 0{,}134 = \mathbf{6{,}0\;\text{mm}}
\end{align*}

\subsection{RTCP --- Remote Tool Center Point (G43.4)}

Sans compensation RTCP, toute rotation de la table déplace le point de contact outil-pièce de manière indésirable. Le RTCP, activé par \code{G43.4}, corrige dynamiquement les axes linéaires lors de toute variation de~A ou~C, maintenant la pointe de l'outil invariante dans le repère pièce.

La figure~\ref{fig:rtcp} illustre le principe~:

\begin{figure}[H]
\centering
\begin{tikzpicture}[scale=0.85]
  % Sans RTCP
  \begin{scope}
    \node[font=\bfseries\small, above] at (2,4.5) {Sans RTCP};
    \draw[thick, gray] (0,0) -- (4,0);
    \fill[orange!30] (1.5,0) rectangle (2.5,0.5);
    \node[font=\tiny] at (2,0.25) {Pièce};
    \draw[very thick, brown!70] (2,3) -- (2,0.7);
    \fill[brown!70] (2,0.7) circle(2pt) node[right, font=\tiny]{TCP\textsubscript{1}};

    % Après rotation
    \draw[thick, gray, dashed] (0,0) -- (4,0);
    \fill[orange!15, rotate around={20:(2,0)}] (1.5,0) rectangle (2.5,0.5);
    \draw[very thick, brown!30] (2,3) -- (2,0.7);
    \fill[red!70] (2.4,0.3) circle(2pt) node[right, font=\tiny, red]{TCP\textsubscript{2} $\neq$ TCP\textsubscript{1} !};
    \draw[-{Stealth}, red, thick] (2,0.7) -- (2.4,0.3);
    \node[font=\tiny\itshape, red] at (2.8,0.8) {Erreur};
  \end{scope}

  % Avec RTCP
  \begin{scope}[xshift=7cm]
    \node[font=\bfseries\small, above] at (2,4.5) {Avec RTCP (G43.4)};
    \draw[thick, gray] (0,0) -- (4,0);
    \fill[orange!30] (1.5,0) rectangle (2.5,0.5);
    \node[font=\tiny] at (2,0.25) {Pièce};
    \draw[very thick, brown!70] (2,3) -- (2,0.7);
    \fill[brown!70] (2,0.7) circle(2pt) node[right, font=\tiny]{TCP};

    % Après rotation + compensation
    \fill[orange!15, rotate around={20:(2,0)}] (1.5,0) rectangle (2.5,0.5);
    \draw[very thick, brown!30] (1.6,3.2) -- (2,0.7);  % outil déplacé
    \fill[green!70] (2,0.7) circle(2pt);
    \draw[-{Stealth}, green!60!black, thick] (2,3) -- (1.6,3.2);
    \node[font=\tiny\itshape, green!60!black] at (1,3.5) {Compensation\\X,Y,Z};
    \node[font=\tiny\bfseries, green!60!black] at (2,0.2) {TCP invariant \checkmark};
  \end{scope}
\end{tikzpicture}
\caption{Principe du RTCP --- sans compensation (gauche) vs. avec G43.4 (droite).}
\label{fig:rtcp}
\end{figure}

L'état RTCP est suivi par \code{isRtcpActive} dans \code{MachineState}, mis à jour par \code{GrblParser.parseModalState()}.

\subsection{Gestion des singularités cinématiques}

La configuration Trunnion présente une singularité à $A = 0°$~: les axes~X et~C deviennent colinéaires (\textit{Gimbal Lock}). La détection logicielle utilise~:

\begin{equation}
|A| < \theta_{sing} = 5° \quad \Rightarrow \quad \text{singularityRisk} = 1 - \frac{|A|}{\theta_{sing}}
\label{eq:singularity}
\end{equation}

Le champ \code{singularityRisk} (0{,}0 $\rightarrow$ 1{,}0) alimente un indicateur visuel dans le dashboard, avertissant l'opérateur de la proximité de la zone singulière. La stratégie de contournement consiste en un lissage par interpolation linéaire pré-calculée dans la zone $|A| < 5°$.

% ══════════════════════════════════════════════════════════════════════
\section{Planification de Trajectoire et Algorithmes Avancés}
% ══════════════════════════════════════════════════════════════════════

\subsection{Streaming Character-Counting}

Le service \code{GCodeStreamingService} implémente l'algorithme \textit{Character-Counting} pour saturer le buffer RX de l'ESP32 sans débordement. La figure~\ref{fig:char_count} illustre le principe.

\begin{figure}[H]
\centering
\begin{tikzpicture}[
  block/.style={draw, rounded corners=3pt, minimum width=2.5cm, minimum height=0.8cm, fill=#1!12, font=\small},
  buf/.style={draw, thick, minimum width=4cm, minimum height=0.8cm, font=\small\ttfamily},
]
  % App side
  \node[block=blue] (app) at (0,0) {Forgeron (IHM)};
  \node[block=blue, below=0.5cm of app] (fg) {\small ForceGuard\\[-2pt]\tiny processLine()};
  \node[block=blue, below=0.5cm of fg] (cc) {\small Character\\[-2pt]\small Counting};

  % Buffer
  \node[buf, right=3cm of cc] (buffer) {Buffer RX};
  \node[font=\scriptsize, above] at (buffer.north) {$B_{max} = 127$ octets};

  % Remplissage visuel du buffer
  \fill[green!30] ([xshift=1pt, yshift=1pt]buffer.south west) rectangle ([xshift=-1.8cm, yshift=-1pt]buffer.north east);
  \node[font=\tiny\bfseries] at ([xshift=-0.7cm]buffer.center) {$B_{vol}$};
  \node[font=\tiny, gray] at ([xshift=1.3cm]buffer.center) {libre};

  % ESP32
  \node[block=red, right=3cm of buffer] (esp) {ESP32\\FluidNC};

  % Flèches
  \draw[-{Stealth}, thick] (app) -- (fg);
  \draw[-{Stealth}, thick] (fg) -- node[left, font=\tiny]{ligne filtrée} (cc);
  \draw[-{Stealth}, thick, blue!60] (cc) -- node[above, font=\tiny]{si $B_{vol} + |L_i| \leq 127$} (buffer);
  \draw[-{Stealth}, thick, green!60!black] (buffer) -- node[above, font=\tiny]{Step/Dir} (esp);
  \draw[-{Stealth}, thick, orange, dashed] (esp.south) -- ++(0,-0.8) -| node[below, font=\tiny, pos=0.25]{\texttt{ok} $\rightarrow$ libère $|L_j|$ octets} (cc.south);

\end{tikzpicture}
\caption{Algorithme Character-Counting --- pipeline de streaming G-code.}
\label{fig:char_count}
\end{figure}

La condition d'envoi est~:
\begin{equation}
\boxed{B_{vol} + |L_i| \leq B_{max} = 127 \;\text{octets}}
\label{eq:char_count}
\end{equation}

Un \textit{watchdog} de 2~secondes surveille les acquittements~: en l'absence de réponse, le streaming est suspendu et l'état machine est interrogé via la commande temps réel \code{?}.

\begin{lstlisting}[caption={Pseudocode --- algorithme Character-Counting.}, label={alg:streaming}]
MAX_RX = 127;  bytes_in_flight = 0
sent_sizes = Queue()

function attemptSend():
    while pendingLines is not empty:
        line = ForceGuard.processLine(pendingLines.first)
        if bytes_in_flight + len(line) <= MAX_RX:
            pendingLines.removeFirst()
            bytes_in_flight += len(line)
            sent_sizes.enqueue(len(line))
            connection.sendRaw(line)
        else: break   // buffer plein --> attendre 'ok'

function handleAck():
    bytes_in_flight -= sent_sizes.dequeue()
    attemptSend()      // relancer
\end{lstlisting}

\subsection{Profil d'accélération S-Curve (Jerk limité)}

Le profil S-Curve décompose chaque segment en \textbf{7~phases cinématiques}~:

\begin{equation}
s(t) = \begin{cases}
\dfrac{1}{6}\,J \cdot t^3 & \text{phase 1 : jerk + constant} \\[4pt]
v_1 t + \dfrac{1}{2}\,a_{max} \cdot t^2 & \text{phase 2 : accél. constante} \\[4pt]
\text{(transition)} & \text{phase 3 : jerk $-$} \\[4pt]
v_{max} \cdot t & \text{phase 4 : vitesse constante} \\[4pt]
\text{(symétrique)} & \text{phases 5--7 : décélération}
\end{cases}
\label{eq:scurve}
\end{equation}

\begin{figure}[H]
\centering
\begin{tikzpicture}[xscale=1.1, yscale=0.8]
  % Axes
  \draw[-{Stealth}] (0,0) -- (10,0) node[right, font=\small]{$t$};
  \draw[-{Stealth}] (0,0) -- (0,5) node[above, font=\small]{$v(t)$};

  % S-Curve
  \draw[very thick, blue!70]
    (0,0) .. controls (0.5,0) and (0.8,1) .. (1.2,1.5)    % Phase 1
    -- (2.5,3)                                               % Phase 2
    .. controls (2.8,3.5) and (3,4) .. (3.5,4)              % Phase 3
    -- (6.5,4)                                               % Phase 4
    .. controls (7,4) and (7.2,3.5) .. (7.5,3)              % Phase 5
    -- (8.8,1.5)                                             % Phase 6
    .. controls (9.2,1) and (9.5,0) .. (10,0);              % Phase 7

  % Trapézoïdal (comparaison)
  \draw[thick, red!50, dashed]
    (0,0) -- (2,4) -- (8,4) -- (10,0);

  % Labels phases
  \foreach \x/\n in {0.6/1, 1.8/2, 3/3, 5/4, 7/5, 8.2/6, 9.5/7} {
    \node[font=\tiny, gray] at (\x, -0.5) {\n};
  }
  \node[font=\tiny\itshape, gray] at (5,-1) {7 phases cinématiques};

  % Légende
  \draw[very thick, blue!70] (6.5,5.5) -- (7.5,5.5) node[right, font=\small]{S-Curve};
  \draw[thick, red!50, dashed] (6.5,5) -- (7.5,5) node[right, font=\small]{Trapézoïdal};

  % v_max
  \draw[thin, dashed, gray] (0,4) -- (10,4);
  \node[font=\scriptsize, gray, left] at (0,4) {$v_{max}$};
\end{tikzpicture}
\caption{Profil S-Curve (7~phases) vs. profil trapézoïdal classique.}
\label{fig:scurve}
\end{figure}

Le planificateur FluidNC maintient un buffer circulaire de 15~blocs (\code{plannerBuffer}) et utilise un algorithme \textit{Look-Ahead} pour pré-calculer les vitesses de jonction entre segments consécutifs.

% ══════════════════════════════════════════════════════════════════════
\section{Développement de l'IHM de Supervision}
% ══════════════════════════════════════════════════════════════════════

\subsection{Choix du framework}

\begin{table}[H]
\centering
\caption{Comparatif des frameworks IHM multi-plateformes.}
\label{tab:framework}
\begin{tabularx}{\textwidth}{l c c c}
\toprule
Critère & Qt/C++ & Electron.js & \textbf{Flutter/Dart} \\
\midrule
Performance native & Excellente & Faible & \textbf{Excellente} \\
Multi-plateforme & Desktop & Desktop + Web & \textbf{Mobile + Desktop + Web} \\
Hot-reload & Non & Partiel & \textbf{Oui ($< 1$~s)} \\
Taille binaire & $\sim$30~Mo & $\sim$150~Mo & \textbf{$\sim$15~Mo} \\
Courbe d'apprentissage & Élevée & Moyenne & \textbf{Faible} \\
Widgets industriels & Qt Widgets & HTML/CSS & \textbf{Material + Custom} \\
\bottomrule
\end{tabularx}
\end{table}

Flutter a été retenu pour sa compilation native ARM/x64, son \textit{hot-reload} sub-seconde et sa portabilité vers Android (tablette d'atelier), Windows (poste de supervision) et Web.

\subsection{Gestion d'état avec Riverpod}

\begin{table}[H]
\centering
\caption{Comparatif des solutions de gestion d'état Flutter.}
\label{tab:state_mgmt}
\begin{tabularx}{\textwidth}{l c c c >{\bfseries}c}
\toprule
Critère & \code{setState} & BLoC & Provider & Riverpod \\
\midrule
Compile-safe & Non & Oui & Non & Oui \\
Testabilité & Faible & Bonne & Moyenne & Excellente \\
Scoping & Widget & Stream & InheritedWidget & Global + scoped \\
Dépendances auto & Non & Non & Limité & Oui \\
\bottomrule
\end{tabularx}
\end{table}

La chaîne réactive est illustrée en figure~\ref{fig:riverpod_chain}.

\begin{figure}[H]
\centering
\begin{tikzpicture}[
  block/.style={rectangle, draw, rounded corners=3pt, minimum width=2.4cm, minimum height=1cm, fill=#1!10, font=\footnotesize\ttfamily, align=center, drop shadow={opacity=0.1}},
  arrow/.style={-{Stealth[length=2.5mm]}, thick, color=MidnightBlue},
]
  \node[block=gray]   (esp)    {ESP32\\FluidNC};
  \node[block=blue, right=0.8cm of esp]   (conn)   {FluidNC\\Connection};
  \node[block=green, right=0.8cm of conn]  (parser) {Grbl\\Parser};
  \node[block=orange, right=0.8cm of parser] (state)  {Machine\\State};
  \node[block=purple, right=0.8cm of state] (stream) {Stream\\Provider};
  \node[block=red, right=0.8cm of stream]  (widget) {Widget\\build()};

  \draw[arrow] (esp)    -- node[above, font=\tiny]{Wi-Fi} (conn);
  \draw[arrow] (conn)   -- node[above, font=\tiny]{raw msg} (parser);
  \draw[arrow] (parser) -- node[above, font=\tiny]{parse()} (state);
  \draw[arrow] (state)  -- node[above, font=\tiny]{Stream} (stream);
  \draw[arrow] (stream) -- node[above, font=\tiny]{watch()} (widget);
\end{tikzpicture}
\caption{Chaîne réactive Riverpod --- du capteur physique au pixel affiché.}
\label{fig:riverpod_chain}
\end{figure}

\subsection{Protocole de communication WebSocket}

La classe \code{FluidNCConnection} (180~lignes) implémente trois mécanismes de résilience~:

\paragraph{Exponential Backoff.}
\begin{equation}
\Delta t_{retry} = \min(2^n, 30) \;\text{secondes}
\label{eq:backoff}
\end{equation}

\paragraph{Heartbeat.} Envoi de \code{?} toutes les 2~s. Timeout de 10~s $\rightarrow$ déconnexion + reconnexion.

\paragraph{Auto-Discovery.} Scan réseau en 3~phases~: IPs favorites ($< 2$~s), sous-réseau /24 (rafales de 20, $\sim$12~s), détection automatique du préfixe.

\begin{figure}[H]
\centering
\begin{tikzpicture}[
  phase/.style={draw, thick, rounded corners=4pt, minimum width=3.5cm, minimum height=1cm, fill=#1!12, font=\small, align=center},
  arrow/.style={-{Stealth[length=3mm]}, very thick, gray!60},
]
  \node[phase=green] (p1) at (0,0) {Phase 1\\{\tiny IPs favorites ($< 2$~s)}};
  \node[phase=blue, right=1.5cm of p1] (p2) {Phase 2\\{\tiny Scan /24 ($\sim$12~s)}};
  \node[phase=purple, right=1.5cm of p2] (p3) {Phase 3\\{\tiny Détection préfixe}};

  \draw[arrow] (p1) -- node[above, font=\tiny\itshape]{trouvé ?} (p2);
  \draw[arrow] (p2) -- node[above, font=\tiny\itshape]{+ résultats} (p3);

  \node[draw, dashed, rounded corners, below=0.8cm of p2, minimum width=4cm, fill=green!5, font=\small] {Appareil trouvé : \code{FluidNC@192.168.1.100}};
\end{tikzpicture}
\caption{Trois phases de l'\code{AutoDiscoveryService}.}
\label{fig:autodiscovery}
\end{figure}

\subsection{Présentation des écrans de l'application}

L'application Forgeron comporte \textbf{six écrans principaux}, conçus avec un layout responsive (mobile, tablette, desktop). Les figures~\ref{fig:screen_dashboard} à~\ref{fig:screen_diagnostics} présentent les captures d'écran de l'application.

\screenshot{Tableau de bord --- DRO 5~axes, statut FSM, overrides, mode 3AX/5AX, température, progression SD}{screen_dashboard}

\screenshot{Palpage \& Origines --- WCS (G54--G59), jogging 5~axes, séquences G38.2 (centre trou, hauteur Z, angle)}{screen_probing}

\screenshot{Magasin d'Outils --- liste des 24~outils, offsets longueur/diamètre, appel M6, durée de vie}{screen_tools}

\screenshot{Espace de Travail --- gestionnaire de fichiers G-code, éditeur, visualisation toolpath}{screen_workspace}

\screenshot{Terminal MDI --- saisie directe G-code, historique, macros rapides (homing, WCS, probe)}{screen_terminal}

\screenshot{Diagnostics --- entrées GPIO, configuration YAML FluidNC, télémétrie réseau, logs}{screen_diagnostics}

% ══════════════════════════════════════════════════════════════════════
\section{Sûreté de Fonctionnement du Code}
% ══════════════════════════════════════════════════════════════════════

\subsection{ForceGuard --- protection des efforts en temps réel}

Le \code{ForceGuardService} est intégré dans le pipeline de streaming. La figure~\ref{fig:forceguard_pipeline} illustre les deux filtres en cascade.

\begin{figure}[H]
\centering
\begin{tikzpicture}[
  block/.style={draw, thick, rounded corners=4pt, minimum width=3.5cm, minimum height=1.2cm, fill=#1!12, font=\small, align=center},
  arrow/.style={-{Stealth[length=3mm]}, thick},
  decision/.style={diamond, draw, thick, fill=yellow!10, font=\small, inner sep=2pt, aspect=2},
]
  \node[block=blue] (input) at (0,0) {Ligne G-code\\entrante};
  \node[decision] (mode) at (4,0) {Mode ?};
  \node[block=red] (filter) at (8,1.5) {\textbf{filterRotaryAxes()}\\{\tiny Bloquer A/C}};
  \node[block=orange] (clamp) at (8,-1.5) {\textbf{clampFeedrate()}\\{\tiny F $\leq$ 500 mm/min}};
  \node[block=green] (output) at (12,0) {Ligne validée\\$\rightarrow$ envoi ESP32};

  \draw[arrow] (input) -- (mode);
  \draw[arrow] (mode) -- node[above left, font=\tiny]{3AX} (filter);
  \draw[arrow] (mode) -- node[below left, font=\tiny]{5AX} (clamp);
  \draw[arrow] (filter) -- (output);
  \draw[arrow] (clamp) -- (output);
\end{tikzpicture}
\caption{Pipeline ForceGuard --- filtrage en cascade selon le mode d'usinage.}
\label{fig:forceguard_pipeline}
\end{figure}

\begin{table}[H]
\centering
\caption{Paramètres ForceGuard par mode d'usinage.}
\label{tab:forceguard}
\begin{tabular}{l c c c c l}
\toprule
Mode & $F_{max}$ & $R_{max}$ & $a_p$ max & $a_e$ max & Action ForceGuard \\
\midrule
\rowcolor{blue!5}
5AX & 500~mm/min & 45,6~N & 0,3~mm & 1,0~mm & Bridage automatique de F \\
\rowcolor{orange!5}
3AX & 2\,000~mm/min & 180~N & 2,0~mm & 6,0~mm & Blocage des axes A/C \\
\bottomrule
\end{tabular}
\end{table}

La méthode \code{checkFile()} permet une vérification pré-usinage complète~: l'intégralité du fichier est analysée et une liste d'alertes \code{ForceGuardAlert} (ligne + message) est présentée à l'opérateur \textbf{avant} le lancement.

\subsection{Intégrité des données et résilience réseau}

La robustesse repose sur quatre mécanismes~:

\begin{enumerate}
  \item \textbf{Typage fort et immuabilité}~: chaque champ de \code{MachineState} est typé et borné (Freezed + \code{@Default}). L'immuabilité (\code{copyWith}) élimine les problèmes de concurrence.
  \item \textbf{Watchdog streaming}~: timeout de 2~s sur les acquittements $\rightarrow$ pause + interrogation \code{?}.
  \item \textbf{Reconnexion automatique}~: exponential backoff (éq.~\ref{eq:backoff}), délai max 30~s.
  \item \textbf{Zombie detection}~: heartbeat 2~s + timeout 10~s $\rightarrow$ déconnexion forcée + reconnexion.
\end{enumerate}

\begin{figure}[H]
\centering
\begin{tikzpicture}[xscale=0.65, yscale=0.7]
  \draw[-{Stealth}] (0,0) -- (16,0) node[right, font=\small]{$t$ (s)};
  \draw[-{Stealth}] (0,0) -- (0,5.5) node[above, font=\small]{État};

  % Niveaux
  \draw[thin, dashed, gray] (0,4.5) -- (15.5,4.5);
  \node[font=\scriptsize, left] at (0,4.5) {Connecté};
  \draw[thin, dashed, gray] (0,1) -- (15.5,1);
  \node[font=\scriptsize, left] at (0,1) {Déconnecté};

  % Timeline
  \fill[green!20] (0,4.2) rectangle (3,4.8);
  \node[font=\tiny] at (1.5,4.5) {Connecté};

  % Perte
  \draw[very thick, red] (3,4.5) -- (3,1);
  \node[font=\tiny\bfseries, red, above] at (3,5) {Perte Wi-Fi};

  % Retries
  \foreach \x/\d in {3/1, 4/2, 6/4, 10/4} {
    \draw[thick, orange] (\x,1) -- (\x,2.5);
    \draw[thick, orange, -{Stealth}] (\x,2.5) -- (\x,1.5);
    \node[font=\tiny, orange] at (\x,3) {\d s};
  }
  \node[font=\tiny\itshape, orange] at (7,3.5) {Exponential backoff};

  % Reconnexion
  \draw[very thick, green!60!black] (14,1) -- (14,4.5);
  \fill[green!20] (14,4.2) rectangle (15.5,4.8);
  \node[font=\tiny\bfseries, green!60!black, above] at (14,5) {Reconnecté};
\end{tikzpicture}
\caption{Mécanisme de reconnexion automatique avec exponential backoff.}
\label{fig:backoff}
\end{figure}

% ══════════════════════════════════════════════════════════════════════
\section{Conclusion}
% ══════════════════════════════════════════════════════════════════════

Ce chapitre a présenté le développement logiciel du système de commande de la CNC 5~axes Trunnion, structuré autour de cinq axes~:

\begin{itemize}
  \item Une \textbf{architecture Clean Architecture} en quatre couches garantissant l'indépendance firmware/IHM et la substitution Mock/FluidNC par inversion de dépendance.
  \item Une \textbf{cinématique inverse analytique} propre au Trunnion, avec compensation RTCP (G43.4) et gestion des singularités ($\Delta\theta = 0{,}019°$).
  \item Un \textbf{streaming industriel} par algorithme Character-Counting (buffer 127~octets, watchdog 2~s) intégrant le ForceGuard pour la protection des efforts en temps réel.
  \item Une \textbf{IHM Flutter/Riverpod} responsive (6~écrans, mobile/desktop) offrant une supervision temps réel (latence $< 5$~ms).
  \item Des mécanismes de \textbf{résilience réseau} (exponential backoff, heartbeat, zombie detection) assurant la continuité de service en environnement industriel Wi-Fi.
\end{itemize}

Le chapitre suivant présentera la réalisation physique du prototype et les essais de validation expérimentale.

% ══════════════════════════════════════════════════════════════════════
\begin{thebibliography}{9}
\bibitem{martin2017clean}
R.~C.~Martin, \textit{Clean Architecture: A Craftsman's Guide to Software Structure and Design}, Prentice Hall, 2017.

\bibitem{grbl2023}
Grbl Contributors, \textit{Grbl v1.1 Interface --- Protocol and Usage}, 2023. [En ligne].

\bibitem{fluidnc2024}
FluidNC Contributors, \textit{FluidNC Documentation --- Configuration and Usage}, 2024. [En ligne].

\bibitem{flutter2024}
Google, \textit{Flutter Documentation --- Build apps for any screen}, 2024. [En ligne].

\bibitem{riverpod2024}
R.~Music, \textit{Riverpod --- A Reactive Caching and Data-binding Framework}, 2024. [En ligne].

\bibitem{iso6983}
ISO~6983-1:2009, \textit{Automation systems and integration --- Numerical control of machines --- Program format and definitions of address words}.
\end{thebibliography}

\end{document}
